\chapter[Implementation and Evaluation]{Implementation and Evaluation of
    SwarmScan}%
\label{chap:implementation}

\section{Introduction}%
\label{sec:impl:intro}

\Cref{chap:system} described the system: what each component does, what it
produces, and why it is built that way. This chapter describes how that design
was actually built, and how it was tested.

The first part covers the implementation. It presents the tools that were used
and the machine everything runs on, then the way the simulator and the flight
software are connected, then the warehouse in which the missions take place, and
finally the concrete realisation of each component, including the values that
\Cref{chap:system} deliberately left aside.

The second part covers the evaluation. It defines the situations the system was
placed in, the reference method it is compared against, the quantities that were
measured, and the results obtained.

\section{Technology Stack and Execution Environment}%
\label{sec:impl:stack}

\subsection{Technologies and Libraries}%
\label{subsec:impl:libraries}

The system runs entirely in simulation. Nothing in it is a model of a drone
written for this work: the vehicle is a physical body inside a physics engine,
and it is flown by the same autopilot software that flies a real aircraft. This
section presents the tools that make that possible, starting from the simulator
and working outwards to the libraries used by each component.

\textbf{NVIDIA Isaac Sim}\footnote{\url{https://developer.nvidia.com/isaac/sim}}
is a robotics simulator built on the Omniverse platform. It provides three things
that this work depends on. It runs the rigid-body physics at a fixed step, so
that a vehicle has real inertia and a command has a real delay. It renders the
camera images with ray tracing, which is what makes a printed
\glsxtrshort{abb:qr} code legible or illegible in an image depending on distance,
angle and lighting. And it simulates a rotating \gls{abb:lidar} that returns a
range for every ray, which is the sensor the shared map is built from. The scene
itself, including the warehouse, the racks and the cartons, is described in the
USD format and manipulated through the \texttt{pxr} library.

\textbf{Pegasus
    Simulator}\footnote{\url{https://pegasussimulator.github.io/PegasusSimulator/}}
is an extension for Isaac Sim that adds aerial vehicles and connects them to real
flight software. It contributes the airframe model---the rotors, the thrust
curve, the drag---and the communication bridge that carries sensor readings out
of the simulator and motor commands back in.

\textbf{The vehicle} is the Iris quadrotor, a four-rotor airframe widely used as a
reference model in drone simulation. Its rotors extend about 0.34\,m from its
centre, which is one of the two terms of the clearance radius used for route
planning in \Cref{subsec:sys:map}, the other being the lateral oscillation
measured while the vehicle holds a position. Its body is a compact shell, and the
two reading cameras are mounted below it so that the hull never appears in a
reading image. The same airframe is declared on both sides of the simulation, in
Pegasus and in the autopilot, so that the vehicle being flown and the vehicle
being simulated have identical characteristics.

\textbf{ArduPilot}\footnote{\url{https://ardupilot.org}} is the autopilot
software. It is used here in \glsxtrshort{abb:sitl} mode, which stands for
\emph{software in the loop}: the complete flight stack runs as an ordinary
program on the same machine, receiving simulated sensor data instead of real ones
and producing motor commands as if it were on board. One such process runs for
each vehicle. This is a deliberate choice rather than a convenience. The system
never writes a position into the simulator; it sends velocity commands to an
autopilot, which decides how to achieve them, and the vehicle therefore behaves
with the delays, the overshoot and the attitude dynamics of a real aircraft. The
autopilots are used in GUIDED mode, in which an external program supplies targets
while the autopilot retains attitude control.

\textbf{MAVLink} is the protocol spoken between the system and the autopilots,
through the
\texttt{pymavlink}\footnote{\url{https://github.com/ArduPilot/pymavlink}}
library. Every take-off, every velocity command, and every position reading
passes through it.

\textbf{NumPy}\footnote{\url{https://numpy.org}} and
\textbf{SciPy}\footnote{\url{https://scipy.org}} carry all the numerical work.
The occupancy grid, the coverage layer, the \gls{abb:lidar} point clouds and the
candidate scoring are all array operations, which is what allows a full
\gls{abb:lidar} revolution to be integrated into the map in a few milliseconds.
SciPy is used in particular for the rotation conversions between the several
frames the system works in.

\textbf{OpenCV}\footnote{\url{https://opencv.org}} handles everything related to
images and camera geometry: colour conversion, the intrinsic matrix, the pose of
a planar square from its four corners, and the drawing of the annotated views
used in the figures and the mission videos. The headless build is used, since no
window is ever opened.

\textbf{zxing-cpp}\footnote{\url{https://github.com/zxing-cpp/zxing-cpp}} is the
\glsxtrshort{abb:qr} decoder retained for flight.
\textbf{ZBar}\footnote{\url{https://github.com/mchehab/zbar}} and
\textbf{BoofCV}\footnote{\url{https://boofcv.org}} were evaluated on the same
images and are kept as alternatives; the comparison that led to this choice is
reported in Section 4.5.2.

\textbf{Ultralytics}\footnote{\url{https://docs.ultralytics.com}} provides the
training and inference framework for the learned detector described in
\Cref{subsec:sys:perception}, in the YOLO11 family. It is used both to train the
detector on the automatically annotated images and to run it during flight.

\textbf{PyTorch}\footnote{\url{https://pytorch.org}} is the underlying
deep-learning framework for both the detector and the vision-language guide.
\textbf{Transformers}\footnote{\url{https://huggingface.co/docs/transformers}}
loads the guide and its processor.
\textbf{bitsandbytes}\footnote{\url{https://github.com/bitsandbytes-foundation/bitsandbytes}}
provides the four-bit quantisation used for the guide.
\textbf{PEFT}\footnote{\url{https://huggingface.co/docs/peft}} applies the
low-rank adapter used by the trained variant of the guide.

\textbf{qrcode}\footnote{\url{https://github.com/lincolnloop/python-qrcode}} and
\textbf{Pillow}\footnote{\url{https://python-pillow.org}} generate the
\glsxtrshort{abb:qr} images that are pasted onto the carton faces when a
warehouse is built, so that the codes read during a mission are genuine
\glsxtrshort{abb:qr} symbols rendered in the scene rather than markers placed by
hand.

\textbf{Matplotlib}\footnote{\url{https://matplotlib.org}} produces the
measurement figures, and \textbf{FFmpeg}\footnote{\url{https://ffmpeg.org}}
encodes the mission videos from the images captured during flight.
\textbf{pytest}\footnote{\url{https://pytest.org}} runs the unit tests that check
each component in isolation, without starting the simulator.

\begin{table}[!htb]
    \centering
    \caption{Technologies and libraries used, with their role in the system.}%
    \label{tab:impl:stack}
    \small
    \setlength{\tabcolsep}{4pt}
    \begin{tabular}{@{}>{\raggedright\arraybackslash}p{3.2cm}
            >{\raggedright\arraybackslash}p{2.1cm}
            >{\raggedright\arraybackslash}p{9.2cm}@{}}
        \toprule
        \textbf{Technology}       & \textbf{Version} & \textbf{Role in the system}                                          \\
        \midrule
        Isaac Sim                 & 5.1.0            & Rigid-body physics, ray-traced rendering, \gls{abb:lidar} simulation \\
        USD (\texttt{pxr})        & ---              & Description of the scene: building, racks, cartons, panels           \\
        Pegasus Simulator         & 5.1              & Airframe model and bridge between the simulator and the flight software \\
        Iris quadrotor            & ---              & The simulated vehicle, four rotors                                   \\
        ArduPilot \glsxtrshort{abb:sitl} & 4.6.0-beta1 & Flight software, one process per vehicle, GUIDED mode              \\
        \texttt{pymavlink}        & 2.4.49           & MAVLink link between the system and the autopilots                   \\
        Python                    & 3.11             & Implementation language                                              \\
        NumPy                     & 1.26.0           & Occupancy grid, coverage layer, point clouds, candidate scoring       \\
        SciPy                     & 1.15.3           & Rotation conversions between the camera, vehicle and world frames     \\
        OpenCV                    & 4.11.0           & Image handling, camera geometry, pose of a planar square              \\
        zxing-cpp                 & 3.1.1            & \glsxtrshort{abb:qr} decoding during flight                          \\
        ZBar (\texttt{pyzbar})    & 0.1.9            & Alternative decoder, evaluated on the same images                     \\
        BoofCV (\texttt{PyBoof})  & 0.43.1           & Alternative decoder, evaluated on the same images                     \\
        Ultralytics               & 8.3.253          & Training and inference of the learned detector, YOLO11 family         \\
        PyTorch                   & 2.7.0            & Deep-learning framework for the detector and the guide                \\
        Transformers              & 5.12.1           & Loading of the vision-language guide and its processor                \\
        bitsandbytes              & 0.50.2           & Four-bit quantisation of the guide                                    \\
        PEFT                      & 0.20.0           & Low-rank adapter for the trained variant of the guide                 \\
        \texttt{qrcode}           & 7.4.2            & Generation of the \glsxtrshort{abb:qr} symbols pasted on the carton faces \\
        Pillow                    & 11.3.0           & Image handling during scene generation                                \\
        Matplotlib                & 3.10.3           & Measurement figures                                                   \\
        FFmpeg                    & 4.4.2            & Encoding of the mission videos                                        \\
        pytest                    & 9.1.1            & Unit tests of each component, without the simulator                   \\
        \bottomrule
    \end{tabular}
\end{table}

\FloatBarrier

Package versions are pinned through a constraints file, since Isaac Sim requires
a specific version of NumPy.

\subsection{Execution Environment}%
\label{subsec:impl:environment}

The simulator, the autopilot processes and the system itself run on a single
machine, whose specifications are given in \Cref{tab:impl:environment}.

\begin{table}[!htb]
    \centering
    \caption{Execution environment.}%
    \label{tab:impl:environment}
    \setlength{\tabcolsep}{4pt}
    \begin{tabular}{@{}>{\raggedright\arraybackslash}p{4.4cm}
            >{\raggedright\arraybackslash}p{10.2cm}@{}}
        \toprule
        \textbf{Component} & \textbf{Specification}                                          \\
        \midrule
        CPU                & Intel Core i9-10900, 10 cores and 20 threads at 2.80\,GHz       \\
        \addlinespace
        RAM                & 31\,GiB                                                         \\
        \addlinespace
        GPU                & NVIDIA GeForce RTX 2060 SUPER, 8\,GB                            \\
        \addlinespace
        GPU driver         & 535.288.01                                                      \\
        \addlinespace
        CUDA               & 12.6                                                            \\
        \addlinespace
        Operating system   & Ubuntu 22.04.5 LTS                                              \\
        \addlinespace
        Kernel             & 6.8.0                                                           \\
        \bottomrule
    \end{tabular}
\end{table}

\FloatBarrier

The graphics memory is shared between the renderer, the learned detector and the
semantic guide. The resulting configuration of each of these components is given
in Sections 4.5.3 and 4.8.
